%% kapa: exam=1 round=23 year=2012 period=1 form=A subject=civil_law subject_ko=민법 seq=6 number=6
%% source: Q-Net 감정평가사 기출문제 (cst00309) artlSeq=5052578 file=raw/1008/5052578/2012년 제23회 감정평가사 제1차시험 기출문제(1교시 A형).pdf
%% conversion: kordoc 4.24.51 -> tools/gen1.py (auto); review=proofread-sonnet; flags=none
\begin{question}{6}
무효와 취소에 관한 설명으로 옳지 않은 것은? (다툼이 있으면 판례에 의함)
\begin{choices}
  \item 의사무능력으로 인한 무효행위의 추인은 그 무효 원인이 소멸한 후에 하여야 효력이 있다.
  \item 취소권 행사기간은 제척기간이므로 당사자의 주장에 관계없이 법원이 당연히 조사하여 고려하여야 한다.
  \item 미성년자는 성년자가 된 후에야 단독으로 추인할 수 있다.
  \item 하나의 계약이라 할지라도 가분성을 가지거나 그 목적물의 일부가 특정될 수 있다면 그 일부만의 취소도 가능하다.
  \item 무효인 법률행위는 추인에 의해 소급하여 유효한 법률행위로 전환된다.
\end{choices}
\end{question}
\answer{5}
